\section{Architecture du pipeline \RAMA{}}
\label{sec:pipeline}

\subsection{Vue d'ensemble}

Le pipeline \RAMA{} comprend cinq étapes séquentielles, chacune implémentée comme un module autonome avec des entrées, sorties et modes de défaillance explicites.

\begin{enumerate}[label=\textbf{Étape~\arabic*:},leftmargin=3cm]
  \item \textbf{Extraction par vision.} Un modèle de vision multimodal (Gemini~2.5 Flash) traite chaque page manuscrite numérisée et retourne un schéma Pydantic structuré contenant~: transcription \LaTeX{} brute, classification d'archétype, score de confiance, et les 12 premiers coefficients de $q$-série.
  \item \textbf{Recherche évolutionnaire génétique.} Un algorithme génétique à base de population recherche l'espace combinatoire des vecteurs d'exposants $\mathbf{r} = (r_1, \ldots, r_{d_{\max}})$ avec $d_{\max} = 12$ et $r_d \in [-24, 24]$.
  \item \textbf{Pont de classification.} L'archétype est utilisé pour assigner un domaine cible et une ombre modulaire conjecturée.
  \item \textbf{Auto-formalisation Lean~4.} Le générateur produit un fichier Lean~4 contenant des théorèmes de Niveau~A (vérification par \texttt{rfl}) et Niveau~B (structure de classification).
  \item \textbf{Cartographie physique.} Les candidats vérifiés sont évalués via une matrice de point-selle.
\end{enumerate}

\subsection{Détails de l'algorithme génétique}

\subsubsection{Fonctionnelle d'énergie}
\begin{equation}
  E(\mathcal{S}) = \alpha \cdot C(\mathcal{S}) + \beta \cdot I(\mathcal{S}) + \gamma \cdot D(\mathcal{S}),
\end{equation}
où $C$ est le coût de complexité (régularisation d'Occam), $I$ l'erreur d'ajustement (distance $\ell_2$ normalisée), et $D$ la pénalité de domaine.

\subsubsection{Opérateurs évolutionnaires}
\begin{itemize}
  \item \textbf{Sélection~:} tournoi avec $k=3$.
  \item \textbf{Croisement~:} uniforme sur les vecteurs d'exposants.
  \item \textbf{Mutation~:} $p_{\mathrm{mut}} = 0{,}3$, $\Delta r_d \sim \mathrm{Uniforme}(-4, 4)$.
  \item \textbf{Immigration~:} injection de 5 individus aléatoires quand la diversité chute.
\end{itemize}

\subsubsection{Hyperparamètres}
\begin{center}
\begin{tabular}{lcc}
\toprule
\textbf{Paramètre} & \textbf{Essai (10 pages)} & \textbf{Exécution complète} \\
\midrule
Taille de population & 25 & 50 \\
Générations max & 5 & 15 \\
$d_{\max}$ & 12 & 12 \\
Taux de mutation & 0,3 & 0,3 \\
Taille de tournoi & 3 & 3 \\
\bottomrule
\end{tabular}
\end{center}

\subsection{Protocole de vérification Lean~4}

\begin{definition}[Certification par le noyau]
Un candidat est \emph{certifié par le noyau} si~: (1)~le fichier Lean~4 compile avec \texttt{lake env lean} retournant le code de sortie~0, (2)~le fichier ne contient aucun \texttt{sorry}, \texttt{axiom}, ni \texttt{native\_decide}, (3)~le script d'audit CI ne rapporte aucun axiome interdit.
\end{definition}

\begin{remark}
La certification par le noyau confirme la \emph{correction au niveau des types} des identités algébriques énoncées. Elle ne vérifie \emph{pas} que le candidat correspond à l'intention originale de Ramanujan ni qu'il est <<~nouveau~>> par rapport à la littérature existante.
\end{remark}
